\chapter{Implementasjon og feilrespons}
\label{ch:implementation}
\section{Motoroppstart og begrenset gjenoppretting}
Et fortsatt spenningssatt PWM-kort kan beholde eldre registerverdier når
micro:biten starter på nytt. Motorpakken setter derfor alle åtte
motorkanalenes FULL\_OFF-bit før oscillatorens oppvåkning. Når
auto-increment er ukjent, adresseres kanalene hver for seg. Etter kjent
oppsett kan en samlet 33-byte-transaksjon slå av alle åtte. Begge varianter
bruker eksplisitt FULL\_OFF etter registerkartet \cite{pca9685_2026}.
Stopp påvirker ikke servoens kanal 14.

I2C-skriving har timeoutparameter 20 ms og rapporterer adresse, register
og bussstatus ved feil. En feil opphever intern motorberedskap og låser
stopp. Initialiseringen kan ikke fullføres delvis og likevel tillate kjøring.
Ingen automatisk busstilbakestilling eller omstart skjer fra feilhandleren.
Hvis bussen er ødelagt, kan heller ikke stopp garanteres via den samme bussen.

\section{Exceptions med tydelig betydning}
\begin{tabularx}{\textwidth}{@{}P{54mm}Y@{}}\toprule
Exception & Betydning\\\midrule
\kode{Motor_Bus_Error} & Motortransaksjonen feilet.\\
\kode{Motor_Not_Ready} & Motor er ikke initialisert eller A har ikke armert.\\
\kode{Servo_Limit_Error} & Vinkelen overskrider tillatt område.\\
\kode{Safety_Setup_Error} & Stoppknappen kunne ikke registreres.\\
\kode{Stop_Latched_Error} & En latchet stopp hindrer videre arbeid.\\
\kode{Lights_Task_Error} & Separat lys-task har publisert feil.\\\bottomrule
\end{tabularx}

Vanlige hindringer og
ugyldige avstandsmålinger er normal kontrollflyt som gir stopp, ikke en
exception for hver måling.

\kode{Adaptive_Fault_Handler} lagrer første feil, setter stoppintensjon og
forsøker én begrenset stoppoperasjon. En sekundær stoppfeil lagres separat,
slik at den ikke skjuler den opprinnelige årsaken. Logging skjer etter
stoppforsøket; en loggerfeil blir registrert og kan ikke starte bilen.
Hovedprogrammet blir deretter i feiltilstanden. Meldingen om akseptert
stopp betyr busskvittering, ikke observert stillstand på hjulene.

Kontrolleren setter sin egen tilstand til Fault før den propagaterer en
exception til denne grensen. Senere kall kan ikke gjenoppta kjøring eller
utløse en storm av nye stoppforsøk. Initialisering er tilsvarende ett forsøk.
Motorfeil, taskfeil og B-stopp krever kontroll av årsaken og en bevisst omstart.

\section{Skjerm og lys uten unødvendig blokkering}
OLED er et tillegg til lokal retningsvisning. Ved første mislykkede
OLED-transaksjon åpner koden en programvarebasert sperre: skjermen
deaktiveres resten av kjøringen, og senere tekster gir ingen busstrafikk.
Dette reduserer både gjentatte ventetider og belastning på den delte bussen.
micro:bitens 5$\times$5-matrise fungerer fortsatt. Uventede Ada-exceptions
skjules ikke av denne degraderingen.

Lys-tasken bruker \kode{delay until} med 50 ms ønsket periode og hopper
over tapte utgivelser. Blinkperioden er ett sekund med 50 prosent påtid.
Et modusskifte starter blinkfasen på, slik at også en kort sving kan vises.
Parkert/brems gir røde lys, kjøring rett frem slukker lysene, sving blinker
på den valgte siden og Fault gir rødt med begge blinklys. Status viser
kommando/intensjon, ikke målt nedbremsing. Ved intern taskfeil publiseres
sticky feil og stoppflagg før tasken parkeres.

\section{Skanning og kontinuerlig fremoverkjøring}
I bevegelse brukes enkeltavlesninger i rekkefølgen front, venstre, front,
høyre. Ved avvisning forsøkes stopp før logging, OLED eller en ny måling.
Et B-trykk under en avlesning kontrolleres før neste handling. Når målingene
er gode, sendes ingen ny start-/stoppkommando for hver runde.

Skanning foregår etter stopp. Frontservoen sentreres, og alle tre sensorer
kontrolleres. Fri front gir null sidepunkter. Ukjent eller utilstrekkelig
svingklarering avslutter søket. Ellers prøves siden med størst sideklaring
først, med forrige svingretning som regel ved likhet. Høyre bruker 72/60 grader,
venstre 96/108 grader. Hvert sidepunkt stabiliserer servoen i 500 ms og
filtrerer tre råprøver. En avvist stråle forkaster kandidatparet straks.

Et bredt nok par krever ny sentrert sensorkontroll. Fersk fri front kan
igjen velge rett frem. Høyst fire sideposisjoner prøves, men sentrerte
kontroller og tre råprøver per punkt kommer i tillegg. Søkets resultat
bevarer tidspunktet til den eldste strålen. Maksimal tillatt alder er
2000 ms, også etter siste live-kontroll rett før svingkommandoen.

Ved sentrering måles de faste sidesensorene mens frontservoen stabiliseres.
Alle tre prøver per sensor beholdes, global triggeravstand gjelder fortsatt,
og frontmålingen venter til minst 500 ms etter sentreringskommandoen.
Det overlapper nyttig arbeid med nødvendig venting uten å innføre flere
sensor-eiere eller korte ned servoventingen.

Den tidligere ekstra servo-/filterrunden i Turn er fjernet: skanningen
returnerer allerede sentrert med en ny kontroll. Turn kontrollerer alle
live-sensorene og tidsgrensen før kjøring, og stopper etter den overvåkede
svingsekvensen. En utløpt kandidat gir stopp. Alderbudsjettet er et
designvalg; tidsbruk og faktisk gjennomførbarhet må testes på målmaskinen.
